\documentclass[11pt]{article}
% Shared preamble for the PNPInverse onboarding docs.
% Each document does % Shared preamble for the PNPInverse onboarding docs.
% Each document does % Shared preamble for the PNPInverse onboarding docs.
% Each document does \input{preamble} after \documentclass.

\usepackage[utf8]{inputenc}
\usepackage[T1]{fontenc}
\usepackage{lmodern}
\usepackage[margin=1in]{geometry}
\usepackage{amsmath, amssymb, amsthm}
\usepackage{mathtools}
\usepackage{empheq}
\usepackage{siunitx}
\usepackage{booktabs}
\usepackage{array}
\usepackage{enumitem}
\usepackage{xcolor}
\usepackage{framed}
\usepackage{microtype}
\usepackage{listings}
\usepackage{tikz}
\usetikzlibrary{arrows.meta, positioning, decorations.pathreplacing, calligraphy}
\usepackage[hidelinks]{hyperref}

% --- branding -------------------------------------------------------------
\definecolor{pnpblue}{RGB}{30,70,120}
\definecolor{pnpgray}{RGB}{90,90,90}
\definecolor{calloutbg}{RGB}{245,245,235}
\definecolor{shadecolor}{RGB}{245,245,235}

% --- callouts (uses framed.sty's shaded environment) ----------------------
\newenvironment{callout}[1]
  {\begin{shaded}\noindent\textbf{\color{pnpblue}#1}\par\smallskip\ignorespaces}
  {\end{shaded}}

\newcommand{\concept}[1]{\textbf{\color{pnpblue}#1}}
\newcommand{\warn}[1]{\textbf{\color{red!70!black}#1}}

% --- code listings --------------------------------------------------------
\lstset{
  basicstyle=\ttfamily\small,
  breaklines=true,
  columns=fullflexible,
  keepspaces=true,
  showstringspaces=false,
  frame=single,
  framesep=4pt,
  backgroundcolor=\color{calloutbg},
  rulecolor=\color{pnpgray!40},
}

% --- math shortcuts -------------------------------------------------------
\newcommand{\R}{\mathbb{R}}
\newcommand{\diff}{\mathrm{d}}
\newcommand{\eqd}{\overset{!}{=}}
\newcommand{\grad}{\nabla}
\newcommand{\divg}{\nabla\!\cdot}
\newcommand{\eqdef}{\;\coloneqq\;}
\newcommand{\evalat}[2]{\left.#1\right\rvert_{#2}}
\newcommand{\set}[1]{\{#1\}}

% common chemistry / physics symbols
\newcommand{\Hyd}{\mathrm{H}}
\newcommand{\Oxy}{\mathrm{O}}
\newcommand{\HtwoO}{\mathrm{H_2O}}
\newcommand{\HtwoOtwo}{\mathrm{H_2O_2}}
\newcommand{\Otwo}{\mathrm{O_2}}
\newcommand{\Hp}{\mathrm{H^+}}
\newcommand{\OHm}{\mathrm{OH^-}}
\newcommand{\Kp}{\mathrm{K^+}}
\newcommand{\Csp}{\mathrm{Cs^+}}
\newcommand{\SOfour}{\mathrm{SO_4^{2-}}}
\newcommand{\Vrhe}{V_{\mathrm{RHE}}}
\newcommand{\eo}{E^{\circ}}
\newcommand{\Eeq}{E_{\mathrm{eq}}}
\newcommand{\kzero}{k_0}

\newcommand{\onboardingfooter}{%
  \vfill\noindent\rule{\linewidth}{0.4pt}\par
  {\small\color{pnpgray} PNPInverse intern onboarding,
  compiled \today. Source under \texttt{docs/onboarding/}.}\par}
 after \documentclass.

\usepackage[utf8]{inputenc}
\usepackage[T1]{fontenc}
\usepackage{lmodern}
\usepackage[margin=1in]{geometry}
\usepackage{amsmath, amssymb, amsthm}
\usepackage{mathtools}
\usepackage{empheq}
\usepackage{siunitx}
\usepackage{booktabs}
\usepackage{array}
\usepackage{enumitem}
\usepackage{xcolor}
\usepackage{framed}
\usepackage{microtype}
\usepackage{listings}
\usepackage{tikz}
\usetikzlibrary{arrows.meta, positioning, decorations.pathreplacing, calligraphy}
\usepackage[hidelinks]{hyperref}

% --- branding -------------------------------------------------------------
\definecolor{pnpblue}{RGB}{30,70,120}
\definecolor{pnpgray}{RGB}{90,90,90}
\definecolor{calloutbg}{RGB}{245,245,235}
\definecolor{shadecolor}{RGB}{245,245,235}

% --- callouts (uses framed.sty's shaded environment) ----------------------
\newenvironment{callout}[1]
  {\begin{shaded}\noindent\textbf{\color{pnpblue}#1}\par\smallskip\ignorespaces}
  {\end{shaded}}

\newcommand{\concept}[1]{\textbf{\color{pnpblue}#1}}
\newcommand{\warn}[1]{\textbf{\color{red!70!black}#1}}

% --- code listings --------------------------------------------------------
\lstset{
  basicstyle=\ttfamily\small,
  breaklines=true,
  columns=fullflexible,
  keepspaces=true,
  showstringspaces=false,
  frame=single,
  framesep=4pt,
  backgroundcolor=\color{calloutbg},
  rulecolor=\color{pnpgray!40},
}

% --- math shortcuts -------------------------------------------------------
\newcommand{\R}{\mathbb{R}}
\newcommand{\diff}{\mathrm{d}}
\newcommand{\eqd}{\overset{!}{=}}
\newcommand{\grad}{\nabla}
\newcommand{\divg}{\nabla\!\cdot}
\newcommand{\eqdef}{\;\coloneqq\;}
\newcommand{\evalat}[2]{\left.#1\right\rvert_{#2}}
\newcommand{\set}[1]{\{#1\}}

% common chemistry / physics symbols
\newcommand{\Hyd}{\mathrm{H}}
\newcommand{\Oxy}{\mathrm{O}}
\newcommand{\HtwoO}{\mathrm{H_2O}}
\newcommand{\HtwoOtwo}{\mathrm{H_2O_2}}
\newcommand{\Otwo}{\mathrm{O_2}}
\newcommand{\Hp}{\mathrm{H^+}}
\newcommand{\OHm}{\mathrm{OH^-}}
\newcommand{\Kp}{\mathrm{K^+}}
\newcommand{\Csp}{\mathrm{Cs^+}}
\newcommand{\SOfour}{\mathrm{SO_4^{2-}}}
\newcommand{\Vrhe}{V_{\mathrm{RHE}}}
\newcommand{\eo}{E^{\circ}}
\newcommand{\Eeq}{E_{\mathrm{eq}}}
\newcommand{\kzero}{k_0}

\newcommand{\onboardingfooter}{%
  \vfill\noindent\rule{\linewidth}{0.4pt}\par
  {\small\color{pnpgray} PNPInverse intern onboarding,
  compiled \today. Source under \texttt{docs/onboarding/}.}\par}
 after \documentclass.

\usepackage[utf8]{inputenc}
\usepackage[T1]{fontenc}
\usepackage{lmodern}
\usepackage[margin=1in]{geometry}
\usepackage{amsmath, amssymb, amsthm}
\usepackage{mathtools}
\usepackage{empheq}
\usepackage{siunitx}
\usepackage{booktabs}
\usepackage{array}
\usepackage{enumitem}
\usepackage{xcolor}
\usepackage{framed}
\usepackage{microtype}
\usepackage{listings}
\usepackage{tikz}
\usetikzlibrary{arrows.meta, positioning, decorations.pathreplacing, calligraphy}
\usepackage[hidelinks]{hyperref}

% --- branding -------------------------------------------------------------
\definecolor{pnpblue}{RGB}{30,70,120}
\definecolor{pnpgray}{RGB}{90,90,90}
\definecolor{calloutbg}{RGB}{245,245,235}
\definecolor{shadecolor}{RGB}{245,245,235}

% --- callouts (uses framed.sty's shaded environment) ----------------------
\newenvironment{callout}[1]
  {\begin{shaded}\noindent\textbf{\color{pnpblue}#1}\par\smallskip\ignorespaces}
  {\end{shaded}}

\newcommand{\concept}[1]{\textbf{\color{pnpblue}#1}}
\newcommand{\warn}[1]{\textbf{\color{red!70!black}#1}}

% --- code listings --------------------------------------------------------
\lstset{
  basicstyle=\ttfamily\small,
  breaklines=true,
  columns=fullflexible,
  keepspaces=true,
  showstringspaces=false,
  frame=single,
  framesep=4pt,
  backgroundcolor=\color{calloutbg},
  rulecolor=\color{pnpgray!40},
}

% --- math shortcuts -------------------------------------------------------
\newcommand{\R}{\mathbb{R}}
\newcommand{\diff}{\mathrm{d}}
\newcommand{\eqd}{\overset{!}{=}}
\newcommand{\grad}{\nabla}
\newcommand{\divg}{\nabla\!\cdot}
\newcommand{\eqdef}{\;\coloneqq\;}
\newcommand{\evalat}[2]{\left.#1\right\rvert_{#2}}
\newcommand{\set}[1]{\{#1\}}

% common chemistry / physics symbols
\newcommand{\Hyd}{\mathrm{H}}
\newcommand{\Oxy}{\mathrm{O}}
\newcommand{\HtwoO}{\mathrm{H_2O}}
\newcommand{\HtwoOtwo}{\mathrm{H_2O_2}}
\newcommand{\Otwo}{\mathrm{O_2}}
\newcommand{\Hp}{\mathrm{H^+}}
\newcommand{\OHm}{\mathrm{OH^-}}
\newcommand{\Kp}{\mathrm{K^+}}
\newcommand{\Csp}{\mathrm{Cs^+}}
\newcommand{\SOfour}{\mathrm{SO_4^{2-}}}
\newcommand{\Vrhe}{V_{\mathrm{RHE}}}
\newcommand{\eo}{E^{\circ}}
\newcommand{\Eeq}{E_{\mathrm{eq}}}
\newcommand{\kzero}{k_0}

\newcommand{\onboardingfooter}{%
  \vfill\noindent\rule{\linewidth}{0.4pt}\par
  {\small\color{pnpgray} PNPInverse intern onboarding,
  compiled \today. Source under \texttt{docs/onboarding/}.}\par}


\title{\bfseries The PNP--BV Mathematical Model \\
       \large from mass conservation to the residual our solver inverts}
\author{Onboarding doc 02 of 06}
\date{\today}

\begin{document}
\maketitle

\begin{callout}{What this document is for}
Builds the equations of the forward solver from scratch.
Assumes you remember what a partial derivative is and have done
multivariable calculus and ODEs; assumes you have \emph{not} taken a
PDE course. By the end you should be able to read the file
\texttt{Forward/bv\_solver/forms\_logc\_muh.py} and recognise every
term as something with a physical name.
\end{callout}

\tableofcontents
\bigskip\hrule\bigskip

% =========================================================================
\section{Why we need PDEs at all}

The simplest model of an electrochemical cell would be a single
\concept{stirred tank}: pretend the electrolyte has uniform
concentrations everywhere, the surface chemistry happens at one
``rate'', and you write down a system of ODEs in time. Many ORR
papers do this.

The problem: the chemistry happens in a region (the electric double
layer + the diffusion layer) that is at most $\sim 100\,\mu$m thick,
inside which concentrations and the electric potential vary by
\emph{many orders of magnitude}. Local pH near the surface can swing
from the bulk value of 4 up to almost 10; the electric potential
$\phi$ can drop by 1 V over a Debye length of less than 1 nm.
Treating any of those as ``uniform'' throws away the physics.

So we keep $c_i(x,t)$ and $\phi(x,t)$ as functions of position
$x \in [0, L_{\mathrm{eff}}]$ and time $t$. Mass conservation and
Gauss's law give us coupled \concept{partial differential
equations}: derivatives in both $x$ and $t$.

\subsection{Steady state}

For analyzing a single $\Vrhe$ point in a current-voltage curve we
hold the applied voltage fixed for long enough that
$\partial_t c_i = 0$. So in practice the residual we solve is the
\emph{steady-state} system. The time-dependent form is in the
\texttt{forms\_logc\_muh.py} \texttt{\_build\_residual} routine but
the time derivative gets multiplied by an artificial relaxation
parameter that we drive to zero (this is sometimes called
``pseudo-time stepping'' and helps Newton's method).

% =========================================================================
\section{Building the PDE system one piece at a time}

\subsection{Mass conservation}

For each chemical species $i$ in the electrolyte, write down its
concentration $c_i(x,t)$ in mol/m$^3$ and its molar flux
$\mathbf{J}_i(x,t)$ in mol/m$^2$/s. Conservation of moles gives
\concept{the continuity equation}:
\begin{equation}
  \frac{\partial c_i}{\partial t} + \divg \mathbf{J}_i = R_i,
  \label{eq:mass}
\end{equation}
where $R_i$ is a volumetric source/sink (units mol/m$^3$/s) due to
homogeneous chemical reactions. For most of our species $R_i = 0$
because we lump all chemistry into a boundary condition at the
electrode. Phase 6$\alpha$ added a non-zero $R_{\Hp}$ from water
self-ionization and Phase 6$\beta$ adds a boundary source from
cation hydrolysis at the OHP.

\subsection{The flux: Nernst--Planck}

The Nernst--Planck equation (doc 01 \S 6.1):
\begin{equation}
  \mathbf{J}_i = -D_i \grad c_i
                 - \frac{z_i F D_i}{RT}\, c_i\, \grad\phi
                 + c_i \mathbf{v}.
  \label{eq:np-2}
\end{equation}
The three terms are diffusion (down the concentration gradient),
migration (charged ions drift along the electric field), and
convection (carried by the bulk velocity field). In our 1D slab the
convection term is folded into the choice of $L_{\mathrm{eff}}$ ---
see doc 01 \S 6.3 --- so the operative flux is just the first two
terms.

\subsection{The potential: Poisson}

In a continuous medium, Gauss's law of electrostatics is
\begin{equation}
  -\divg(\varepsilon\, \grad\phi) = \rho,
  \label{eq:poisson}
\end{equation}
where $\varepsilon$ is the dielectric permittivity of the solvent
(for water at 25$^{\circ}$C, $\varepsilon = 80\,\varepsilon_0$) and
$\rho$ is the local volumetric charge density. The charge density is
just the sum of ion charges:
\begin{equation}
  \rho = F \sum_i z_i\, c_i.
  \label{eq:rho}
\end{equation}
(\ref{eq:poisson}) and (\ref{eq:rho}) close the system: $\phi$ is
determined by the ion distribution, and the ion distribution is
shaped by $\phi$ through (\ref{eq:np-2}). This nonlinear feedback is
the heart of EDL physics.

\subsection{Putting it together: the PNP system}

\begin{empheq}[box=\fbox]{align}
  \frac{\partial c_i}{\partial t}
     + \divg\!\Big(-D_i \grad c_i - \tfrac{z_i F D_i}{RT}\,c_i\,\grad\phi\Big) &= R_i
     \quad \text{for each species } i \\
  -\divg(\varepsilon\, \grad\phi) &= F \sum_i z_i\, c_i
\end{empheq}
This is a coupled, nonlinear, parabolic--elliptic system. In 1D
steady state on $x \in (0, L_{\mathrm{eff}})$ it becomes a
coupled boundary-value problem.

% =========================================================================
\section{Boundary conditions}

\subsection{At the bulk side $x = L_{\mathrm{eff}}$ (easy)}

Here we are far enough from the electrode that the EDL has fully
relaxed and concentrations recover their bulk values. Dirichlet:
\[
  c_i(L_{\mathrm{eff}}) = c_i^{\mathrm{bulk}}, \qquad
  \phi(L_{\mathrm{eff}}) = 0.
\]
The choice of $\phi=0$ at the bulk is just a gauge: we measure
electric potential relative to the bulk solution.

\subsection{At the electrode side $x = 0$ (the interesting one)}

Three conditions hold simultaneously here:

\paragraph{(a) Faradaic flux from BV.} The flux of each species
into the boundary equals its stoichiometric draw in the surface
reactions. For a reaction $\sum_p \nu_p \,\mathrm{S}_p + n_e e^- \to
\sum_q \nu_q \,\mathrm{S}_q$ proceeding at rate $r$
[mol/m$^2$/s]:
\begin{equation}
  \mathbf{J}_i \cdot \hat{\mathbf{n}}\big|_{x=0}
     = \sum_{\text{rxns}} \nu_i \, r,
  \label{eq:bv-flux}
\end{equation}
where $\hat{\mathbf{n}}$ is the outward normal (pointing into the
bulk). $r$ is the Butler--Volmer rate (the BV equation in doc 01,
or its log-rate form). For us, summing over R$_{2e}$
and R$_{4e}$ in parallel, the boundary flux of $\Otwo$ is
$\nu_{\Otwo}^{(2e)} r_{2e} + \nu_{\Otwo}^{(4e)} r_{4e}$ with
$\nu_{\Otwo} = -1$ for both reactions.

\paragraph{(b) The Stern compact layer.} The metal-side surface
charge $\sigma_{\mathrm{metal}}$ is set by the rigid Stern
capacitor as
\begin{equation}
  \sigma_{\mathrm{metal}} = C_S \,(\phi_m - \phi(0)) = C_S \,\psi_S,
  \label{eq:stern}
\end{equation}
where $\phi_m$ is the metal potential, $\phi(0)$ is the diffuse-side
solution potential at the OHP, and $\psi_S$ is the
\concept{Stern voltage drop}. Continuity of the electric
displacement gives a Robin condition for $\phi$ at $x=0$:
\begin{equation}
  -\varepsilon \frac{\partial \phi}{\partial x}\bigg|_{x=0}
     = -\sigma_{\mathrm{metal}}
     = -C_S\,(\phi_m - \phi(0)).
  \label{eq:phi-bc}
\end{equation}
The applied potential $V$ enters through $\phi_m$ via
$\phi_m = V - V^{\circ}_{\mathrm{ref}}$ for whatever reference is
chosen.

\paragraph{(c) Equilibrium for non-reactive species.} For ions that
do not react at the surface (e.g.\ $\SOfour$, $\Kp$ in our
baseline), the flux into the boundary is zero:
$\mathbf{J}_i\cdot \hat{\mathbf{n}}|_{x=0} = 0$. In our
\concept{analytic Bikerman counter-ion} treatment we go further: we
\emph{don't even keep} a NP equation for these ions, but instead
substitute their equilibrium Boltzmann profile (Bikerman form, doc 01)
directly. They act as a closed-form contribution to the
charge density in (\ref{eq:rho}). This saves a degree of freedom per
species and is exact at steady state for non-reactive ions.

% =========================================================================
\section{The Bikerman closure for multi-ion electrolytes}

For non-reactive analytic counter-ions we use a
\concept{multi-ion shared-$\theta$ Bikerman closure}. In the
single-ion form it reduced to the standard Bikerman expression (doc 01). The multi-ion form
that the codebase implements (in
\texttt{Forward/bv\_solver/boltzmann.py}) is:
\begin{equation}
  c_i(x) = \frac{c_i^{\mathrm{bulk}}\,\exp(-z_i \phi(x)/V_T)\,(1 - A_{\mathrm{dyn}}(x))}
                {\theta_b + \sum_j a_j\, c_j^{\mathrm{bulk}}\,\exp(-z_j \phi(x)/V_T)},
  \label{eq:multi-bikerman}
\end{equation}
where $a_i$ is the molar volume of species $i$, $\theta_b$ is the
bulk packing fraction, and $A_{\mathrm{dyn}}(x)$ is the local
volume fraction occupied by the dynamic species (the ones for which
we still solve a NP equation). This makes the steric saturation a
\emph{shared} budget: the dynamic species crowding takes up part of
the available volume; the analytic Bikerman ions split what's left
according to their bulk concentrations and their Boltzmann factors.

The single hard rule about this:
\begin{callout}{Hard rule (don't violate)}
The IC seed and the residual must agree about steric saturation.
The initializer (\texttt{set\_initial\_conditions\_debye\_boltzmann\_*})
seeds the composite-$\psi$ matched-asymptotic with multi-species
$\gamma$; the residual side calls
\texttt{build\_steric\_boltzmann\_expressions}. Mixing a Bikerman IC
with an ideal-counter-ion residual (or vice versa) cold-fails on
the saturated manifold. This is documented as ``Hard rule \#3'' in
the project \texttt{CLAUDE.md}.
\end{callout}

% =========================================================================
\section{Choosing the primary variables wisely}

\subsection{Why log-concentration ($u_i = \ln c_i$)}

Concentrations in our problem range from $10^{-12}$ mol/m$^3$
(deep-cathodic $\Hp$ depletion) up to $\sim 10^{4}$ mol/m$^3$
(close-packed counter-ion) --- 16 orders of magnitude. Solving for
$c_i$ directly causes Newton iterates to step into negative
concentrations and crash. Instead we change variable to
\begin{equation}
  u_i = \ln c_i, \qquad c_i = e^{u_i},
  \label{eq:u}
\end{equation}
which automatically keeps $c_i > 0$ no matter what Newton step is
taken. The diffusion and migration terms (\ref{eq:np-2}) become:
\[
  \mathbf{J}_i = -D_i\, c_i \,\grad u_i
                 - \tfrac{z_i F D_i}{RT}\, c_i\, \grad\phi
               = -D_i\, c_i\, \grad\!\left(u_i + \tfrac{z_i F}{RT}\phi\right).
\]
Notice the bracketed combination $u_i + (z_i F/RT)\phi$ is exactly
the dimensionless electrochemical potential
$\tilde\mu_i / RT$ minus its standard-state piece. This is the cue
for the next step.

\subsection{Why $\mu$ for protons (the \texttt{logc\_muh}
            formulation)}

Even with $u_i = \ln c_i$, the proton variable $u_H$ and the
potential $\phi$ each vary by tens of $V_T$ across the diffuse
layer, while their \emph{sum} (the electrochemical potential $\mu_H$)
varies much more gently. The combination
\[
  \mu_H = u_H + e_m\, z_H\, \phi/V_T \quad (z_H = +1)
\]
(where $e_m$ is the dimensionless electromigration prefactor, equal
to 1 in physical units) becomes our new primary variable. Then $u_H$
is recovered as $u_H = \mu_H - z_H \phi/V_T$. Newton iterates on
$\mu_H$ are smooth where iterates on $u_H$ would oscillate.

This is the meaning of the formulation tag \texttt{formulation =
"logc\_muh"}. The other species ($\Otwo$, $\HtwoOtwo$) stay in
plain $u_i = \ln c_i$ since they are uncharged ($z = 0$) and their
behaviour is dominated by diffusion, not migration.

\paragraph{Practical bookkeeping.} The dispatcher in
\texttt{Forward/bv\_solver/dispatch.py} routes the residual builder
to either \texttt{forms\_logc.py} (all species in $\ln c$) or
\texttt{forms\_logc\_muh.py} (proton in $\mu_H$, others in $\ln c$).
Production work is on \texttt{forms\_logc\_muh.py}; the legacy
\texttt{forms\_logc.py} is kept for cross-check.

% =========================================================================
\section{Nondimensionalization}

The solver works in scaled (``hatted'') variables to keep matrix
conditioning sane. The scales are:
\begin{center}
\begin{tabular}{ll}
\toprule
\textbf{Quantity} & \textbf{Scale} \\
\midrule
position $x$    & $\hat{x} = x/L_{\mathrm{eff}}$ \\
concentration $c_i$ & $\hat{c}_i = c_i / c_{\mathrm{ref}}$ where $c_{\mathrm{ref}} = $ bulk $\Otwo$ typically\\
potential $\phi$    & $\hat{\phi} = \phi / V_T$ \\
time $t$        & $\hat{t} = t \cdot D_{\mathrm{ref}} / L_{\mathrm{eff}}^2$ \\
rate $r$        & $\hat{r} = r \cdot L_{\mathrm{eff}} / (D_{\mathrm{ref}} c_{\mathrm{ref}})$ \\
$\kzero$        & rate constant scaled the same way as $r$ \\
\bottomrule
\end{tabular}
\end{center}

Two dimensionless numbers fall out of the scaling:
\begin{itemize}[leftmargin=*]
\item The \concept{electromigration prefactor} $e_m = 1$ (it goes
      away in the natural scaling but is kept as a knob for
      diagnostic use).
\item The \concept{Debye number}
      $\hat{\lambda}_D^2 = \varepsilon V_T / (F^2 c_{\mathrm{ref}}
      L_{\mathrm{eff}}^2)$, which is tiny ($\sim 10^{-10}$). It
      multiplies the Poisson Laplacian, so the EDL is a
      \emph{singular perturbation} layer of width
      $\hat{\lambda}_D \sim 10^{-5}$ in scaled units. The mesh has
      to resolve this layer; doc 03 explains how.
\end{itemize}

The constants live in \texttt{Nondim/} (scaling transforms +
physical constants) and \texttt{scripts/\_bv\_common.py} (production
constants like \texttt{V\_T}, \texttt{C\_O2\_BULK\_HAT},
\texttt{K0\_HAT\_R2E}, etc.). Cross-check any dimensional formula
by computing in SI and then dividing by the scales above.

% =========================================================================
\section{Phase 6$\alpha$ and 6$\beta$ extensions}

The above is the production stack as of Phase 5$\gamma$. Two
extensions are layered on top.

\subsection{Phase 6$\alpha$: water self-ionization}

Adds the homogeneous reaction $\HtwoO \rightleftharpoons
\Hp + \OHm$ with equilibrium $K_w = c_{\Hp}\, c_{\OHm}$. Because
$\OHm$ is not a tracked species, we close it algebraically:
$c_{\OHm}(x) = K_w / c_{\Hp}(x)$. The $\Hp$ source becomes
\[
  R_{\Hp}(x) = k_w^{\mathrm{eff}} (1 - c_{\Hp}(x) c_{\OHm}(x) / K_w)
\]
which vanishes at equilibrium. The flag is
\texttt{enable\_water\_ionization=True}; default is off and the
default-off code path is byte-equivalent to pre-6$\alpha$.

\subsection{Phase 6$\beta$: cation hydrolysis at the OHP}

Adds a coverage variable $\Gamma_{\mathrm{MOH}}(t)$ at the OHP
satisfying:
\begin{equation}
  \dot\Gamma_{\mathrm{MOH}}
     = k_{\mathrm{hyd}}\,(1 - \Gamma_{\mathrm{MOH}}/\Gamma_{\max})\,c_M(0)
     - k_{\mathrm{des}}\,\Gamma_{\mathrm{MOH}},
  \label{eq:gamma}
\end{equation}
where:
\begin{itemize}[leftmargin=*]
\item $c_M(0)$ is the cation concentration at the OHP from the
      Bikerman closure;
\item $k_{\mathrm{hyd}}$ depends on the local field through Singh's
      field-dependent pK$_a$ (Singh 2016 formula, doc 01);
\item the Langmuir factor $(1 - \Gamma/\Gamma_{\max})$ caps the
      coverage at one monolayer; this was added in v10a after v9
      converged at $\sim$6 monolayers (unphysical);
\item $k_{\mathrm{des}}$ is the desorption rate (the inverse step).
\end{itemize}
Each forward step of $\Gamma_{\mathrm{MOH}}$ injects a proton flux
into the boundary condition for $\Hp$, modifying the BV reaction
rates indirectly. The whole thing is solved as an outer Picard loop
around the inner Newton solve for the PNP+BV residual.

The relevant code is in
\texttt{Forward/bv\_solver/cation\_hydrolysis.py}. Calibration of
$\Gamma_{\max}$, $k_{\mathrm{hyd}}$, $k_{\mathrm{des}}$ to literature
values is the v10b work documented in
\texttt{docs/phase6/v10b\_calibration\_summary.md}.

% =========================================================================
\section{From PDE to a residual: a quick FEM tour}

Doc 03 covers this in detail; here is the quick version so you can
read \texttt{forms\_logc\_muh.py} without panic.

\subsection{Weak form}

You cannot solve PDEs exactly on a computer. The
\concept{finite-element method (FEM)} chooses a finite-dimensional
function space $V_h$ (e.g.\ piecewise-linear functions on a mesh) and
demands that the residual of the PDE be \emph{orthogonal to} every
test function $v \in V_h$:
\[
  \mathcal R(c_h, \phi_h)[v] = 0 \qquad \forall v \in V_h.
\]
For a NP equation like $\partial_t c + \divg \mathbf{J} = R$ on the
domain $\Omega$, you multiply by a test function $v$, integrate
over $\Omega$, integrate the divergence by parts, and pull out
boundary terms:
\[
  \int_\Omega \partial_t c \, v \,\diff x
  - \int_\Omega \mathbf{J}\cdot \grad v \,\diff x
  + \int_{\partial\Omega} (\mathbf{J}\cdot \hat{\mathbf{n}})\, v \,\diff s
  - \int_\Omega R\, v\, \diff x = 0.
\]
The boundary integral is exactly where the Butler--Volmer flux
(\ref{eq:bv-flux}) and the Stern term (\ref{eq:phi-bc}) enter.
Firedrake (the FEM library we use) lets you write each of these
integrals as one Python expression on a UFL form
(\texttt{F = inner(J, grad(v))*dx + bv\_flux*v*ds(electrode)}, and
so on); see \texttt{forms\_logc\_muh.py} lines $\sim$200--400.

\subsection{Newton}

After discretization, the residual becomes a vector function $F(U)$
of the global degrees-of-freedom vector $U$ (one entry per mesh node
per primary variable). We solve $F(U) = 0$ by Newton's method:
\[
  J(U^k) \,\Delta U = -F(U^k), \qquad U^{k+1} = U^k + \Delta U,
\]
with $J = \partial F/\partial U$ assembled symbolically by Firedrake
and the linear solve done by PETSc. You have seen Newton in
numerical methods; the only twist here is that $J$ is a sparse
matrix of size $\sim$$10^4$--$10^5$, and that $F$ is so nonlinear
that Newton from cold-start frequently fails. Doc 03 covers the
\concept{continuation} tricks that get around this.

% =========================================================================
\section{What goes in, what comes out}

\paragraph{Inputs to the solver.}
\begin{itemize}[leftmargin=*]
\item Geometry: $L_{\mathrm{eff}}$, mesh $N_x \times N_y$, grading
      factor for the EDL refinement.
\item Bulk concentrations $c_i^{\mathrm{bulk}}$.
\item Reaction set: $\{E^{\circ}_r, k_0^{(r)}, \alpha^{(r)}, n_e^{(r)}\}$ for
      $r \in \{2e, 4e\}$.
\item Stern capacitance $C_S$.
\item Counterion list (Bikerman entries with charge, molar volume,
      and bulk concentration).
\item Applied voltage $\Vrhe$ (often a list, swept).
\end{itemize}

\paragraph{Outputs.}
\begin{itemize}[leftmargin=*]
\item Per voltage point: convergence flag, total disk current $i_D$,
      peroxide current $i_P$, surface concentrations
      $c_i(0)$, surface electric potential $\phi(0)$, $\Gamma_{\mathrm{MOH}}$
      coverage, BV-rate diagnostics.
\item A JSON dump under \texttt{StudyResults/<study\_name>/}.
\item Plots, when the driver script generates them.
\end{itemize}

\onboardingfooter
\end{document}
